\chapter{Documentare și stadiul actual}
\etapa{S1-E2}

În cadrul acestei etape inițiale, documentarea se axează pe înțelegerea conceptelor fundamentale de Asigurare a Calității (QA) în mediul profesional, pe stadiul actual al testării software (manuale și automate) și pe stabilirea modului în care procesele moderne integrează instrumente bazate pe inteligență artificială.

\section{Surse și soluții existente}
Testarea software reprezintă un pilon esențial în dezvoltarea aplicațiilor moderne, asigurând livrarea unor produse stabile și conforme cu cerințele utilizatorilor. Conform literaturii de specialitate și bunelor practici din industrie (precum standardele ISTQB), procesul de QA combină:
\begin{itemize}
    \item Testarea manuală exploratorie și bazată pe scenarii/user stories.
    \item Testarea automatizată a interfețelor (UI) și a serviciilor web (API) pentru creșterea eficienței.
    \item Utilizarea instrumentelor moderne și a tehnicilor de prompt engineering asistate de AI pentru generarea datelor de test și optimizarea cazurilor de testare.
\end{itemize}

\section{Planul individual de lucru}
Planul de practică se desfășoară pe parcursul celor două semestre (112 ore total, împărțite în 14 săptămâni pe semestru) și este structurat pe următoarele etape:
\begin{itemize}
    \item \textbf{E1:} Tema, traseul, depozitul de lucru, convenția și stabilirea regulilor de acces.
    \item \textbf{E2:} Documentarea, fundamentarea testării și explorarea inițială a aplicației.
    \item \textbf{E3:} Analiza cerințelor și proiectarea soluției (șabloane de teste și raportarea defectelor).
    \item \textbf{E4:} Primele componente funcționale (scrierea seturilor de teste și execuția lor).
    \item \textbf{E5:} Integrarea, datele de test și structura caietului.
    \item \textbf{E6:} Versiunea de etapă testată și testele de regresie.
    \item \textbf{E7:} Caietul semestrului și pregătirea prezentării (demo).
\end{itemize}